% =============================================================================
% Conclusion section for MADONNA paper (200–300 words, UK English).
% Paste the content below into \section{Conclusions} in your main .tex file.
% =============================================================================

\section{Conclusions}

This paper presented MADONNA, a static magnetic anomaly detection and analysis testbed comprising two spatially separated observatories (OBS-1 and OBS-2), each equipped with tri-axial magnetometers, an IMU, and a GNSS module. The hardware and data pipeline were described from sensor readout through to central storage and retrieval, supporting long-duration, time-stamped magnetic field recordings for both offline analysis and real-time monitoring. The system uses a GRU-based predictor with cyclical time-of-day features to learn the ambient baseline and an adaptive threshold on prediction residuals to flag anomalies. Detected anomalies are characterised by the direction of the magnetic perturbation vector, giving azimuth and inclination relative to the observatory, and visualised as a unit vector in a three-dimensional plot. Where both observatories record an anomaly within a short time window, triangulation can be used to estimate an approximate source location.

The framework provides a practical foundation for studying magnetic anomalies in an open environment: it separates natural diurnal variation from short-lived disturbances, supports cross-sensor and cross-site correlation, and excludes flagged anomalies from retraining so that the learned baseline remains representative of the quiet field. The direction-finding method correctly computes the direction of the perturbation vector from the three component readings and a quiet-period baseline; however, in field experiments the reported direction can disagree with known source positions. Possible causes include the fact that the direction of the perturbation need not coincide with the direction to the source, baseline choice and drift, sensor alignment and coordinate frame, and the time alignment of the three components when forming the perturbation vector. Addressing these inconsistencies—for example through improved baseline selection, sensor calibration, tighter time alignment of components, and comparison with alternative methods such as dipole fitting or multi-point inversion—remains an important direction for future work. Extending the testbed in these ways would strengthen the reliability of direction finding and support the longer-term goals of directional source estimation and autonomous navigation.
